\section{The induction}\label{sec:induction}

The preceding lemmas will be applied on a succession of scales. We
first select a prefix whose total piece order is controlled by one
of its layers. A fixed window beyond that layer makes the surviving
edges suitable for batching and folding.

Fix the absolute parameters
\begin{equation}\label{eq:induction-parameters}
 \theta=0.95,\qquad K=200,\qquad
 P_0=4(K+1),\qquad \eta=\frac1{100P_0}.
\end{equation}
In particular,
\begin{equation}\label{eq:scale-window}
 1.02\theta^K<0.001,\qquad \theta^{K+1}<0.001,
\end{equation}
since \(\theta^K\le e^{-10}\).

\begin{lemma}[Selection of a prefix]\label{lem:prefix-choice}
Let \((S_j)_{j\ge0}\) be a nonnegative, finitely supported sequence
which is not identically zero. There is an index \(i\) with
\(S_i>0\) such that
\begin{equation}\label{eq:prefix-charge}
 \sum_{j=0}^{i+K}S_j\le P_0S_i.
\end{equation}
\end{lemma}

\begin{proof}
Start at the first positive term and maintain
\(\sum_{j=0}^iS_j\le2(K+1)S_i\). If
\eqref{eq:prefix-charge} fails, the sum of the next \(K\) terms
is greater than \(2(K+1)S_i\). Their largest value \(M\) is
therefore greater than \(2S_i\). Move to an index \(i'\) in
that window with \(S_{i'}=M\). Then
\[
 \sum_{j=0}^{i'}S_j
 \le 2(K+1)S_i+KM
 <(2K+1)M\le2(K+1)S_{i'}.
\]
The invariant persists. Each move increases the index and lands
at a positive term. Finite support therefore forces the procedure
to stop, at which point \eqref{eq:prefix-charge} holds.
\end{proof}

For real \(x\ge1\), define the nondecreasing function
\begin{equation}\label{eq:induction-potential}
 \phi(1)=\frac12,\qquad
 \phi(x)=\max\left\{\frac12,
                1-\frac1{\sqrt{\log x}}\right\}
 \quad(x>1).
\end{equation}
The small improvement over a linear bound will absorb the overlap
introduced when boxes are split.

\begin{proposition}\label{prop:induction-bound}
There is an absolute constant \(C\) such that every finite simple
graph on \(n\ge1\) vertices has an edge partition into at most
\(C\phi(n)n\) simple cycles and single edges.
\end{proposition}

\begin{proof}
We use strong induction on \(n\). The lower scale cutoff \(D_*\)
is an absolute constant, chosen to satisfy the uniform requirements
below. Only after fixing it will we choose \(C\). All auxiliary
constants introduced before that final choice are independent of
\(C\) and remain valid if the cutoff is increased.

Edgeless graphs have the empty partition. For \(1\le n<D_*\),
using every edge singly proves the assertion when \(C\ge D_*\),
since \(\phi(n)\ge1/2\). Now let \(n\ge D_*\), assume the
assertion for every smaller positive order, and put
\[
 D_j=n^{\theta^j},\qquad
 m=\min\{j\ge0:D_j<D_*\}.
\]
Take \(D_*>2\). Then \(m\) is finite and positive, and the
scales \(D_0,\ldots,D_{m-1}\) all satisfy the cutoff.

\paragraph{Prospective layers.}
Starting with the given graph \(G\), perform both operations of
\Cref{lem:scale-split} at each successive scale
\(D_0,\ldots,D_{m-1}\). The inputs at stage \(j\) have average
degree at most \(D_j\): this holds initially because
\(\overline d(G)\le n-1<D_0\), and the residual graph in each
new box has average degree at most
\(D_j^{0.95}=D_{j+1}\).

Let \(\mathcal Y_j\) be the indexed family of level-\(j\)
boxes and \(\mathcal R_j\) the indexed family of pieces removed
inside them. The residual graph passed to the next stage retains
the entire vertex set of its box, including isolated vertices.
Every box has a unique indexed parent, and its vertex set is
contained in its parent's vertex set. Descendants of any box
cover that box's vertex set. Distinct indices are retained even
when their vertex sets overlap or coincide.

These operations define prospective layers: we will subsequently
keep only a prefix of them. At the end of every prefix, the long
cycles removed so far, the pieces removed so far, and the remaining
input graphs partition \(E(G)\). In particular, discarding the
prospective operations beyond a chosen prefix means retaining
their input graphs, with all their assigned edges, at that
prefix's boundary.

Write
\[
 N_{-1}=n,\qquad N_j=\sum_{Y\in\mathcal Y_j}|Y|,
 \qquad S_j=\sum_{R\in\mathcal R_j}|V(R)|,
\]
and set \(S_j=0\) for \(j\ge m\). Increase the absolute
constant \(C_0\) in \Cref{lem:scale-split} if necessary so that
\(C_0\ge1\), and put \(C_2=40C_0\). For \(0\le j<m\),
\[
 \sum_{l=0}^j\frac1{\log D_l}
 =\frac1{\log D_j}\sum_{t=0}^j\theta^t
 \le\frac{20}{\log D_j}.
\]
Provided \(\log D_*\ge C_2\), the product of the order
increases is consequently bounded by
\begin{equation}\label{eq:layer-order}
 n\le N_j
 \le n\prod_{l=0}^j\left(1+\frac{C_0}{\log D_l}\right)
 \le n\left(1+\frac{C_2}{\log D_j}\right)
 \le2n.
\end{equation}
Here we used \(1+x\le e^x\) and \(e^x\le1+2x\) for
\(0\le x\le1/2\). The lower bound follows from preservation
of the vertex cover at every split. The same product calculation,
using the relevant subproduct, shows that the orders of all
level-\(j\) descendants of a box of order \(y\) sum to at most
\((1+C_2/\log D_j)y\). Their union still covers its full
vertex set.

We also require \(D^{-0.01}\le1/\log D\) for every
\(D\ge D_*\). If \(\mathcal L_j\) is the family of long
cycles removed through level \(j\), then
\begin{equation}\label{eq:prefix-cycle-cost}
 |\mathcal L_j|
 \le\sum_{l=0}^j\frac{N_{l-1}}{2D_l^{0.01}}
 \le n\sum_{l=0}^j\frac1{\log D_l}
 \le\frac{20n}{\log D_j}.
\end{equation}
These bounds have no dependence on the number of layers.
Piece disjointness within each box gives
\(S_j\le N_j\le2n\).

\paragraph{Terminal cases.}
If every \(S_j\) vanishes, keep all prospective stages. Use the
long cycles and take the remaining edges singly. The remaining
graphs, on the level-\((m-1)\) boxes, have average degree at most
\(D_m<D_*\), so their total number of edges is at most
\(N_{m-1}D_m/2<D_*n\). By
\eqref{eq:prefix-cycle-cost}, the total cost is at most
\((20+D_*)n\).

Otherwise choose \(i\) by \Cref{lem:prefix-choice}. If
\(i+K\ge m-1\), again keep every stage. Retain all long cycles as
completed parts and take every edge of the final boundary graphs
singly, at cost at most \((20+D_*)n\) as above. The zero extension of
\((S_j)\) ensures that
\(\sum_{j=0}^{m-1}S_j\le P_0S_i\). Apply
\Cref{lem:small-residue} to each piece at its own scale, with
\(c=1\) and the fixed \(\eta\). This completes at most
\(9P_0S_i\) parts and leaves graphs of total order at most
\[
 \eta P_0S_i=\frac{S_i}{100}\le\frac n{50}<\frac n4.
\]
Each individual residual graph has order at most \(\eta r<n\),
where \(r\) is the order of its parent piece: the piece's
vertex set is a subset of \(V(G)\), so \(r\le n\), and
\(\eta<1\). Apply the induction hypothesis to each nonempty
residual graph; their total additional cost is at most \(Cn/4\),
using \(\phi\le1\).

Both terminal cases are therefore covered by the bound
\[
 T(D_*)n+\frac{Cn}{4},\qquad
 T(D_*)=20+D_*+18P_0.
\]
For \(C\ge4T(D_*)\), this is at most
\(Cn/2\le C\phi(n)n\). This choice also covers the base cases.

\paragraph{A prefix with a scale gap.}
It remains to consider \(i+K<m-1\). Keep the stages through
\(t=i+K\), and put
\begin{equation}\label{eq:batch-parameters}
 \begin{gathered}
 D=D_i,\qquad w=\log D,\qquad H=C_2\theta^{-K},\\
 B_0=D^{0.30},\qquad a=D^{0.01},\qquad
 L_0=D^{0.03},\qquad \sigma=D^{0.90}.
 \end{gathered}
\end{equation}
Since \(\log D_t=\theta^K w\),
\eqref{eq:layer-order} and \eqref{eq:prefix-cycle-cost} imply
\begin{equation}\label{eq:gap-global-bounds}
 N_t\le(1+H/w)n,\qquad |\mathcal L_t|\le Hn/w.
\end{equation}
The latter uses \(C_2\ge40>20\).

Fix a level-\(i\) box \(Y\), of order \(y\le D^{1.02}\).
Let \(T_Y\) be the sum of the orders of its level-\(t\)
descendants. Each such descendant carries its remaining graph
after stage \(t\); it has order at most
\(D^{1.02\theta^K}<D^{0.001}\) and that graph has average
degree at most \(D^{\theta^{K+1}}\). Moreover,
\begin{equation}\label{eq:descendant-order}
 y\le T_Y\le(1+H/w)y\le2y.
\end{equation}
For the last inequality, note that
\(H/w=C_2/\log D_t\le1\).

Let \(F(Y)\) be the graph on \(Y\) whose edges are exactly
the remaining edges in these descendants. Their assigned edge
sets are disjoint, so \(F(Y)\) is simple and
\begin{equation}\label{eq:descendant-degree}
 2|E(F(Y))|\le T_YD^{\theta^{K+1}}
 \le2yD^{\theta^{K+1}}.
\end{equation}

Each descendant has fewer than \(L_0\) vertices, so we first group
descendants to make larger folding sets possible. The number of
batches will bound the loss when an intersection of a batch with a
selected piece contains fewer than \(L_0\) vertices of degree at
most \(a\) in \(F(Y)\). The batch order cap \(B_0\) keeps quotient
cycles within the routing team bound.

Partition the indexed descendants into batches whose sums of
orders are at most \(B_0\). This can be done with at most
\(1+4y/B_0\) batches: require the cutoff to ensure
\(D^{0.001}\le B_0/2\), fill a batch until the next descendant
would exceed \(B_0\), and then start a new one. Every batch
except possibly the last has order sum greater than \(B_0/2\),
so their number is at most \(1+2T_Y/B_0\).

For each batch let \(W\) be the union of its vertex sets, and
let \(J_W\) contain exactly its remaining assigned edges. The
batch index is retained in this notation. Denote the family by
\(\mathcal W(Y)\). Thus the \(J_W\) partition \(E(F(Y))\),
the sets \(W\) cover \(Y\), and
\begin{equation}\label{eq:batch-orders}
 |W|\le B_0,\qquad
 \sum_{W\in\mathcal W(Y)}|W|\le T_Y,\qquad
 \sum_{Y\in\mathcal Y_i}\sum_{W\in\mathcal W(Y)}|W|
 \le N_t.
\end{equation}
The order sums used in forming batches include all descendant
multiplicities; passing to a union within a batch can only
decrease them.

Since each level-\(t\) box has a unique indexed level-\(i\)
ancestor, the batch edge sets over all \(Y\) partition the boundary
remainder. Thus, at the boundary of the committed prefix, the
assigned edge sets give the disjoint union
\begin{equation}\label{eq:global-edge-ledger}
 E(G)=
 \mathop{\dot\bigcup}_{L\in\mathcal L_t}E(L)
 \;\dot\cup\;
 \mathop{\dot\bigcup}_{j=0}^{t}\ 
             \mathop{\dot\bigcup}_{R\in\mathcal R_j}E(R)
 \;\dot\cup\;
 \mathop{\dot\bigcup}_{Y\in\mathcal Y_i}\ 
             \mathop{\dot\bigcup}_{W\in\mathcal W(Y)}E(J_W).
\end{equation}

\paragraph{Active sets and multiplicities.}
Call a vertex of \(Y\) good if its degree in \(F(Y)\) is at
most \(a\). By \eqref{eq:descendant-degree} and
\eqref{eq:scale-window}, the number of other vertices is at most
\[
 2yD^{\theta^{K+1}-0.01}\le2yD^{-0.009}\le y/w,
\]
where the final inequality is another uniform large-scale
requirement. For each \(W\in\mathcal W(Y)\) and each
level-\(i\) piece \(R\) belonging to \(Y\), consider all
good vertices of \(W\cap V(R)\). If their number is at least
\(L_0\), retain this set as \(X_{W,R}\) and call it active;
otherwise skip it. Active sets in a fixed batch are disjoint,
because the pieces in \(Y\) have disjoint vertex sets.

There are at most \(y/\sigma\) such pieces. Consequently the
total size of all skipped sets in \(Y\), counting occurrences
in different batches separately, is at most
\begin{align}
 L_0(1+4y/B_0)y/\sigma
 &\le y\bigl(D^{-0.87}+4D^{-0.15}\bigr)
 \le y/w.\label{eq:skipped-occurrences}
\end{align}
Here the exponent \(-0.15\) uses \(y\le D^{1.02}\), and the
last inequality is imposed for all scales above the cutoff.

Let \(s(Y)\) be the sum of the orders of the level-\(i\)
pieces in \(Y\), and define
\[
 \Sigma_Y=\sum_{\substack{W\in\mathcal W(Y),\ R\text{ in }Y\\
                         X_{W,R}\text{ active}}}|X_{W,R}|,
 \qquad \Sigma=\sum_{Y\in\mathcal Y_i}\Sigma_Y.
\]
Every good piece vertex occurs in at least one batch, since the
batches cover \(Y\). Choose one such occurrence per vertex.
At most \(y/w\) piece vertices are not good, and at most
\(y/w\) of the chosen occurrences can lie in skipped sets by
\eqref{eq:skipped-occurrences}. Hence
\(\Sigma_Y\ge s(Y)-2y/w\).

For the upper bound, let \(u_Y(v)\ge1\) be the number of
batches in \(Y\) containing \(v\). Then
\[
 \sum_{v\in Y}\bigl(u_Y(v)-1\bigr)
 =\sum_{W\in\mathcal W(Y)}|W|-y.
\]
Since the pieces in \(Y\) are vertex-disjoint, their excess
occurrences are bounded by this full excess. Thus
\(\Sigma_Y\le s(Y)+\sum_W|W|-y\). Summing the two bounds,
using \(N_i\le2n\), \(N_i\ge n\), and
\eqref{eq:gap-global-bounds}--\eqref{eq:batch-orders}, gives
\begin{equation}\label{eq:active-order}
 S_i-\frac{4n}{w}\le\Sigma
 \le S_i+N_t-N_i\le S_i+\frac{Hn}{w}.
\end{equation}
This counts a nongood vertex only once per indexed box in the
lower bound; no bound on its number of batch occurrences is
needed. Overlaps between different boxes are already included
in both \(S_i\) and \(N_i\).

\paragraph{Resolving the batches and the pieces.}
In every level-\(i\) piece \(R\), apply
\Cref{lem:edge-split} with two classes to partition its assigned
edges into spanning graphs \(P_R\) and \(T_R\), each with
cut expansion at least \(\sigma/4\). Reserve \(P_R\) as a
router and leave \(T_R\) untouched during the batch work.

All hypotheses of \Cref{lem:batch-resolution} now hold in each
\(Y\). The \(J_W\) have order at most \(B_0\) and disjoint
assigned edges; their union is \(F(Y)\). The pieces have
orders between \(\sigma\) and \(D^{1.02}\), disjoint vertex
sets within \(Y\), and edges disjoint from \(F(Y)\).
Every vertex of an active \(X_{W,R}\) has
degree at most \(a\) in the whole \(F(Y)\), and its active
set has order at least \(L_0\). In particular the degree bound
is shared across all batches, as required by that lemma.

The lemma produces simple quotient graphs \(Q_W\), with orders
\begin{equation}\label{eq:individual-quotient-order}
 q_W=|W|-\sum_{R:X_{W,R}\text{ active}}
                       \floor{|X_{W,R}|/2}
 \le |W|-\frac13\sum_{R:X_{W,R}\text{ active}}|X_{W,R}|.
\end{equation}
Every batch has positive order and pair identifications leave
at least one vertex, so \(1\le q_W\le B_0\). Also
\(B_0=D^{0.30}\le n^{0.30}<n\). Apply the induction
hypothesis to these quotients, retaining isolated quotient
vertices in their order counts. The batch lemma lifts all the
resulting quotient parts simultaneously. In each \(Y\), the
additional number of completed parts is at most \(8\Sigma_Y\),
and the remaining graph on every \(R\) still contains \(T_R\).
In particular, lifting a quotient cycle produces one simple
cycle and has no additional part charge.

Apply \Cref{lem:small-residue}, with the fixed \(\eta\), to
the remaining graph of every level-\(i\) piece using \(c=1/4\).
Apply it to every other piece of the committed prefix at its
own scale using \(c=1\). The total number of parts supplied
by these applications is at most \(9P_0S_i\). If
\(\mathcal H\) is their family of residual graphs, then
\begin{equation}\label{eq:piece-residual-order}
 \sum_{H'\in\mathcal H}|V(H')|
 \le\eta\sum_{j=0}^tS_j\le\eta P_0S_i.
\end{equation}
Every nonempty member has order at most \(\eta r<n\), for
the same reason as in the terminal case. Induction applies to
all of them.

Within each \(Y\), \Cref{lem:batch-resolution} resolves exactly
the batch category in \eqref{eq:global-edge-ledger} together
with the router edges it consumes.
Every consumed router edge comes from exactly one level-\(i\)
piece and is used once. The piece lemmas are applied only after
those edges have been removed, and partition all remaining
piece edges into completed parts and the members of
\(\mathcal H\). Distinct indexed boxes have disjoint assigned
edge sets even when their vertex sets overlap. Thus the only
unresolved original edges after lifting the quotient partitions
are precisely the assigned edges of \(\mathcal H\), each
appearing once. Recursively partitioning them completes an edge
partition of \(G\).

\paragraph{Order savings and the final estimate.}
Define the absolute constants
\[
 A_1=H+\frac43,\qquad
 A_2=\max\{9P_0+8,\,9H\}.
\]
By \eqref{eq:individual-quotient-order},
\eqref{eq:batch-orders}, and \eqref{eq:active-order},
\begin{equation}\label{eq:total-quotient-order}
 \sum_{Y,W}q_W\le N_t-\frac{\Sigma}{3}
 \le n-\frac{S_i}{3}+\frac{A_1n}{w}.
\end{equation}
The parts completed apart from the quotient partitions and
the recursive work on \(\mathcal H\) number at most
\begin{align}
 |\mathcal L_t|+8\Sigma+9P_0S_i
 &\le (9P_0+8)S_i+9Hn/w\notag\\
 &\le A_2(S_i+n/w).\label{eq:nonrecursive-cost}
\end{align}
Monotonicity of \(\phi\), \(q_W\le B_0\), and
\eqref{eq:piece-residual-order} therefore bound the complete
partition size by
\begin{equation}\label{eq:inductive-cost}
 A_2(S_i+n/w)+C\eta P_0S_i
 +C\phi(B_0)\left(n-\frac{S_i}{3}+\frac{A_1n}{w}\right).
\end{equation}
The coefficient of \(S_i\) in this expression is at most
\[
 A_2+\frac C{100}-\frac C6
 =A_2-\frac{47C}{300}\le0
 \qquad\text{if }C\ge7A_2.
\]
It remains to pay for the order-overlap term. Set
\(\delta=1/\sqrt{0.30}-1>0\). For \(w\ge40/3\), both
\(\phi(n)\) and \(\phi(B_0)\) take their nonconstant
values, and \(\log n\ge w\) gives
\begin{equation}\label{eq:potential-gain}
 \phi(n)-\phi(B_0)
 =\frac1{\sqrt{0.30w}}-\frac1{\sqrt{\log n}}
 \ge\frac{\delta}{\sqrt w}.
\end{equation}
For \(C\ge7A_2\), the remaining overhead in
\eqref{eq:inductive-cost}, beyond \(C\phi(B_0)n\), is at
most \(C(A_1+1)n/w\). It is absorbed by
\eqref{eq:potential-gain} whenever
\begin{equation}\label{eq:cutoff-absorption}
 \log D_*\ge\left(\frac{A_1+1}{\delta}\right)^2.
\end{equation}
No choice made here requires the total quotient order in
\eqref{eq:total-quotient-order} to be at most \(n\); the
potential gain pays for its possible excess.

We finally specify the order of constant choices. The parameters
in \eqref{eq:induction-parameters}, the scale-split constant
\(C_0\ge1\), and hence \(C_2,H,A_1,A_2,\delta\), are
already absolute and fixed. Choose \(D_*\) large enough for
\Cref{lem:scale-split,lem:batch-resolution}, for
\Cref{lem:edge-split} with two classes at expansion
\(D^{0.90}\), and for \Cref{lem:small-residue} with the
fixed \(\eta\) and \(c\in\{1,1/4\}\). Also impose
\(\log D_*\ge\max\{C_2,40/3\}\),
\eqref{eq:cutoff-absorption}, and, for every \(D\ge D_*\),
\[
 \begin{gathered}
 D^{-0.01}\le\frac1{\log D},\qquad
 D^{0.001}\le\tfrac12D^{0.30},\\
 2D^{-0.009}\le\frac1{\log D},\qquad
 D^{-0.87}+4D^{-0.15}\le\frac1{\log D}.
 \end{gathered}
\]
These are uniform eventual inequalities between fixed powers
and logarithms. Such a cutoff exists independently of \(C\).
After fixing it, choose
\[
 C\ge\max\{7A_2,\,4T(D_*)\}.
\]
The base and terminal cases then hold, and
\eqref{eq:inductive-cost}--\eqref{eq:potential-gain} give at
most \(C\phi(n)n\) parts in the remaining case. This closes
the strong induction.
\end{proof}

Since \(\phi(n)\le1\), \Cref{prop:induction-bound} proves
\Cref{thm:main} for positive order. The zero-vertex graph has
the empty partition, completing all cases.
